% ==========================================================================
%  Section 5 -- The carrier and its Borel structure.  Owner: W05.
%  Frozen labels: sec:carrier, def:carrier-hminus3, prop:completion,
%  lem:coordinate-embedding, thm:separable, thm:cylinder-borel,
%  lem:lift-representative, rem:lift-development, rem:topology.
%
%  Nothing probabilistic is proved in this file: no moment estimate, no
%  tightness, no Prokhorov argument (those are sec:tightness).  What is proved
%  is that the carrier is a complete separable metric space, that the
%  coordinate maps are continuous and separate points, and that the sigma
%  algebra they generate is the Borel sigma algebra.  The open physical-lattice
%  input is stated as a hypothesis (H) and marked as used nowhere in the chain.
% ==========================================================================
\section{The carrier and its Borel structure}\label{sec:carrier}

This section builds the space on which the lattice laws of \cref{sec:lattice}
are Borel probability measures, so that convergence below is asserted against
bounded continuous functions on a concrete space. That space is an explicit
weighted sequence space over the mode lattice, and the properties a limit
theorem normally assumes of its ambient space are proved for it here:
completeness and separability (\cref{def:carrier-hminus3,thm:separable}),
continuity and point separation of the coordinate maps
(\cref{lem:coordinate-embedding}), and the equality of the sigma algebra
generated by finitely many coordinates with the Borel sigma algebra
(\cref{thm:cylinder-borel}). \Cref{prop:completion} records the comparison of
the two weights that makes the compactness argument of \cref{sec:tightness}
possible.

No moment is estimated and no tightness is proved here; that is
\cref{sec:tightness}. Nor is \cref{lem:lift-representative} proved here: it is
stated below and proved from facts of \cref{sec:lattice} and of this section.

\subsection{The carrier}

\begin{definition}[The carrier]\label{def:carrier-hminus3}
The \emph{carrier} is the weighted sequence space
\begin{equation}\label{eq:carrier}
  \mathcal C:=\HL{3}
  =\Bigl\{c=(c_k)_{k\in\Z^2}\in\C^{\Z^2}:
  \sum_{k\in\Z^2}w_3(k)\abs{c_k}^2<\infty\Bigr\},
  \qquad
  \norm{c}_{\mathcal C}^2:=\sum_{k\in\Z^2}w_3(k)\abs{c_k}^2 ,
\end{equation}
with $w_3$ the weight \eqref{eq:weights} of \cref{def:modes}, that is
$w_3(0)=1$ and $w_3(k)=\norm{k}_\infty^{-6}$ for $k\neq0$. The space $\HL{3}$
of \cref{def:carrier-hs} is spelled $\mathcal C$ from here on; it is the space
the measures below live on, and $\norm{\cdot}_{\mathcal C}$ is the norm in which
they are tight. Writing $\widehat c:=(\sqrt{w_3(k)}\,c_k)_{k\in\Z^2}$ for the
\emph{weighted family} of $c$, the map
$\Theta:\mathcal C\to\ell^2(\Z^2)$, $\Theta c=\widehat c$, is a linear isometry
of $\mathcal C$ onto the unweighted space of complex square-summable families on
the mode lattice: the carrier is $\ell^2$ of the mode lattice in weighted
coordinates. The weighted family is what the carrier stores --- the completion
of \cref{def:carrier-hs} is realised as an $\ell^2$ space whose coordinate at
$k$ is $\sqrt{w_3(k)}\,c_k$ --- so below the \emph{coordinates} of a carrier
point are the entries of $\widehat c$, while $c$ is the family of Fourier
coefficients. Its index set is $\Z^2$ throughout: the carrier, its weight and
(\cref{rem:dimension}) the measure chain are two-dimensional, and no
general-dimensional carrier is defined here.
\end{definition}

\begin{remark}[Which family a formula is written in]\label{rem:two-families}
Two families denote the same carrier point and the paper uses both, so the rule
is stated once. The \emph{coefficient} family $c$ is what the symbols and
pairings of \cref{sec:covariance} are written for, and its norm is
\eqref{eq:carrier}. The \emph{stored} family $\widehat c=\Theta c$ is what a
formula written in coordinates uses --- the Fourier lift of
\cref{def:finite-source-lift}, the moment of \cref{sec:tightness}, and the
cylinder maps of \cref{sec:marginals} --- and every weight those formulae carry is
the one transported along $\Theta$, so their coordinates already include it.
Because $\Theta$ is an isometry, a computation may be done in either family; the
only error available is to apply the weight twice, and no formula below does:
a factor $\sqrt{w_3(k)}$ appears in a coordinate formula exactly when the
coefficient at $k$ is being turned into a coordinate, and nowhere else.
\end{remark}

\begin{remark}[One spelling for three source notations]\label{rem:carrier-notations}
The sources write this one space in three ways: as
$\lmod{lp}\ ((\texttt{Fin 2}\to\Z)\to\C)\ 2$, in its curried form
$\lmod{lp}\ (\texttt{Fin 2}\to\Z\to\C)\ 2$, and as the same $\lmod{lp}$
space over a constant family of fibres, the fibres being indexed by the type that
the sources call \lmod{PeriodicFourierIndex}. The second is the curried form of
the first, the third differs only in presenting the family of fibres as a
function of the index, and all three denote the same weighted $\ell^2$ space; the
machine development names it \lmod{PeriodicHMinusThreeCompletion}. This paper
writes $\mathcal C$ and nothing else.
\end{remark}

\begin{proposition}[Completion and embedding]\label{prop:completion}
\mbox{}
\begin{enumerate}
\item $\mathcal C$ is complete. Equivalently, the completion of the finitely
  supported families for the norm $\norm{\cdot}_{\mathcal C}$ is $\mathcal C$
  itself, and $\Theta$ identifies it isometrically with $\ell^2(\Z^2)$.
\item The weights are pointwise ordered, $w_3(k)\le\widetilde w_2(k)$ for every
  $k\in\Z^2$, and consequently
  \begin{equation}\label{eq:embedding}
    \norm{c}_{\mathcal C}\le\norm{c}_{\HL{2}}
    \qquad\text{for every } c\in\HL{2},
  \end{equation}
  so $\HL{2}\subseteq\mathcal C$ and the inclusion is continuous with operator
  norm one, attained at every mode with $\norm{k}_\infty=1$. Moreover
  $w_3(k)/\widetilde w_2(k)=\norm{k}_\infty^{-2}$ tends to $0$ along the mode
  lattice.
\item Let $S\subseteq\HL{2}$ with $\sup_{c\in S}\norm{c}_{\HL{2}}\le B<\infty$.
  Then for every integer $R\ge0$ the members of $S$ have uniformly small tails
  in the carrier norm:
  \begin{equation}\label{eq:tail-estimate}
    \sup_{c\in S}\ \sum_{\norm{k}_\infty>R}w_3(k)\abs{c_k}^2
    \;\le\;\frac{B^2}{(R+1)^2}\ .
  \end{equation}
\end{enumerate}
The compactness statement corresponding to \eqref{eq:tail-estimate}, that the
closure of such a set is compact in $\mathcal C$, is \cref{thm:compact-tail}.
\end{proposition}

\begin{proof}
(i) $\Theta$ is a linear isometry of $\mathcal C$ onto the complete space
$\ell^2(\Z^2)$, and a linear isometry onto a complete space has complete domain,
since Cauchy sequences and limits transport along an isometry.

(ii) For $k\neq0$ we have $\norm{k}_\infty\ge1$, hence
$\norm{k}_\infty^{-6}\le\norm{k}_\infty^{-4}$; at $k=0$ both weights equal $1$.
Summing against $\abs{c_k}^2$ gives \eqref{eq:embedding}. Since
$\norm{e_k}_{\mathcal C}=\norm{e_k}_{\HL{2}}=1$ at every mode of
$\ell^\infty$-norm one, the operator norm is exactly one.

(iii) For integer $R$ and $\norm{k}_\infty>R$ we have $\norm{k}_\infty\ge R+1$,
so $w_3(k)\le(R+1)^{-2}\widetilde w_2(k)$ there. Hence
\[
  \sum_{\norm{k}_\infty>R}w_3(k)\abs{c_k}^2
  \le\frac{1}{(R+1)^2}\sum_{k\in\Z^2}\widetilde w_2(k)\abs{c_k}^2
  =\frac{\norm{c}_{\HL{2}}^2}{(R+1)^2}
  \le\frac{B^2}{(R+1)^2},
\]
the middle identity being the definition of the $\HL{2}$ norm, which uses the
full weight and so includes the zero mode.
\end{proof}

The pair $(2,3)$ is fixed by two different requirements in dimension two.
Gaussian second-moment summability in $\HL{s}$ needs $s>(d-2)/2$, which in $d=2$
asks only for $s>0$, whereas relative compactness of the corresponding sublevel
sets needs a strictly weaker ambient exponent $r>s$; the pair $(s,r)=(2,3)$
satisfies both with room to spare. Only this two-dimensional instance is used
below.

\subsection{Coordinates, separability, and cylinders}

\begin{lemma}[Coordinate embeddings]\label{lem:coordinate-embedding}
\mbox{}
\begin{enumerate}
\item For every $k\in\Z^2$ the coordinate functional
  $\pi_k:\mathcal C\to\C$, $\pi_k(c)=c_k$, is continuous, with
  $\abs{\pi_k(c)}\le w_3(k)^{-1/2}\norm{c}_{\mathcal C}$.
\item For every $m\ge1$, every $m$-tuple of modes and every choice of a real or
  an imaginary part at each slot, the map $\mathcal C\to\R^m$ whose components
  are those parts of $c_{k_1},\dots,c_{k_m}$ is continuous, and so is the
  complex version $\mathcal C\to\C^m$, $c\mapsto(c_{k_1},\dots,c_{k_m})$.
\item The coordinates separate points: if $c\neq c'$ in $\mathcal C$ then
  $c_k\neq c'_k$ for some $k\in\Z^2$.
\end{enumerate}
\end{lemma}

\begin{proof}
(i) $w_3(k)\abs{c_k}^2$ is a single term of the sum \eqref{eq:carrier}, and
$w_3(k)>0$ at every mode. (ii) Each component composes the continuous
functional $\pi_{k_i}$ with one of the continuous maps $z\mapsto\Real z$,
$z\mapsto\Im z$, and a map into a finite product is continuous exactly when its
components are. (iii) If $c_k=c'_k$ for every $k$ the two families coincide.
\end{proof}

Point separation is weaker than determination of the Borel structure, and the
difference matters. Continuity of the coordinates makes the sigma algebra they
generate a sub-sigma-algebra of the Borel sigma algebra, and separation of points
makes the coordinate map injective, but neither says that this
sub-sigma-algebra is the whole Borel sigma algebra: for a metric space that is
not separable the conclusion can fail. The next two results supply separability
first and the sigma algebra equality second.

\begin{theorem}[Separability]\label{thm:separable}
The carrier contains a countable dense subset. Consequently it is a complete
separable metric space, hence Polish, so Prokhorov's theorem applies to families
of Borel probability measures on it (\cref{ext:prokhorov}).
\end{theorem}

\begin{proof}
Fix an enumeration $k_1,k_2,\dots$ of $\Z^2$, and for $m\ge1$ let
$T_mc:=\sum_{j\le m}c_{k_j}e_{k_j}$ keep the first $m$ coordinates in that
enumeration and set the rest to zero. Let $D$ be the set of families
$\sum_{j\le m}q_je_{k_j}$ with $m\ge1$ and $q_j\in\mathbb{Q}+i\mathbb{Q}$; as a
countable union of countable sets, $D$ is countable.

\emph{The truncations converge in the carrier norm.} Let $c\in\mathcal C$. By
\eqref{eq:carrier} the series $\sum_{j\ge1}w_3(k_j)\abs{c_{k_j}}^2$ is
summable, hence converges as a series of non-negative terms, and therefore
\begin{equation}\label{eq:truncation-tail}
  \norm{c-T_mc}_{\mathcal C}^2
  =\sum_{j>m}w_3(k_j)\abs{c_{k_j}}^2\;\longrightarrow\;0
  \qquad(m\to\infty).
\end{equation}
This is where summability enters, and it is the only place: $c-T_mc$ is the tail
of that same series, its squared carrier norm is exactly that tail, and the tail
of a convergent series of non-negative terms tends to zero. What is summed is
therefore the weight against one square-summable family, and
\eqref{eq:truncation-tail} is convergence in the carrier norm.

\emph{Rational approximation at fixed truncation.} Fix $m$ and $\varepsilon>0$.
Since the sum in \eqref{eq:truncation-tail} is over the finite set
$k_1,\dots,k_m$, choose $q_1,\dots,q_m\in\mathbb{Q}+i\mathbb{Q}$ with
$w_3(k_j)\abs{c_{k_j}-q_j}^2\le\varepsilon^2/m$ for $j\le m$; then
$d:=\sum_{j\le m}q_je_{k_j}\in D$ satisfies
$\norm{T_mc-d}_{\mathcal C}\le\varepsilon$.

Given $c$ and $\varepsilon>0$, choose $m$ with
$\norm{c-T_mc}_{\mathcal C}<\varepsilon$ and then $d\in D$ with
$\norm{T_mc-d}_{\mathcal C}<\varepsilon$, so $\norm{c-d}_{\mathcal C}<2\varepsilon$.
Hence $D$ is dense and $\mathcal C$ is separable. Completeness is
\cref{prop:completion}(i), and a complete separable metric space is Polish. The
machine development proves this statement as
\lean{periodicHMinusThreeCompletion_separable}, one of the results listed in
\cref{app:ledger}.
\end{proof}

\begin{theorem}[Cylinder sigma algebra equals Borel]\label{thm:cylinder-borel}
Let $\Sigma_{\mathrm{cyl}}$ be the sigma algebra on $\mathcal C$ generated by the
coordinate functionals $\pi_k$, $k\in\Z^2$: the smallest sigma algebra
containing $\Pi^{-1}(B)$ for every finite coordinate map $\Pi$ of
\cref{lem:coordinate-embedding} and every Borel $B\subseteq\C^m$. Then
$\Sigma_{\mathrm{cyl}}=\mathcal B(\mathcal C)$.
\end{theorem}

\begin{proof}
Write $E:\mathcal C\to\C^{\Z^2}$, $E(c)=(c_k)_{k\in\Z^2}$, for the coordinate
map, and let $\mathcal B_\Pi$ be the product sigma algebra on $\C^{\Z^2}$. Since
$\Z^2$ is countable, $\mathcal B_\Pi$ is generated by the coordinate projections,
so
\begin{equation}\label{eq:cylinder-comap}
  \Sigma_{\mathrm{cyl}}=E^{-1}(\mathcal B_\Pi)
\end{equation}
by definition.

\emph{The cylinders are Borel.} Each $\pi_k$ is continuous by
\cref{lem:coordinate-embedding}(i), hence Borel, so $\pi_k^{-1}(B)$ is Borel for
Borel $B$; the sets with Borel preimage under $E$ form a sigma algebra, and
$\Sigma_{\mathrm{cyl}}\subseteq\mathcal B(\mathcal C)$ follows from
\eqref{eq:cylinder-comap}.

\emph{The Borel sets are cylindrical.} By \cref{thm:separable} the carrier is
separable, so its topology has a countable base: with a countable dense set
$D_0$ fixed, the balls $B(q,r)$, $q\in D_0$, $r\in\mathbb{Q}$, $r>0$, form one.
A sigma algebra containing every base element contains every open set, hence
every Borel set. So it suffices to place every ball in $\Sigma_{\mathrm{cyl}}$.
Fix $c_0\in\mathcal C$, $r>0$ and put
$s_m(c):=\sum_{j\le m}w_3(k_j)\abs{c_{k_j}-c_{0,k_j}}^2$. Each $s_m$ is a finite
sum of functions of one coordinate, hence is
$\Sigma_{\mathrm{cyl}}$-measurable, and $s_m(c)\uparrow\norm{c-c_0}_{\mathcal
C}^2$ because the series \eqref{eq:carrier} converges. Therefore
\[
  \overline{B(c_0,r)}
  =\bigcap_{n\ge1}\bigcup_{M\ge1}\bigcap_{m\ge M}
  \bigl\{c\in\mathcal C:\ s_m(c)\le(r+1/n)^2\bigr\}
\]
is a countable Boolean combination of $\Sigma_{\mathrm{cyl}}$-measurable sets,
and the equality of sets follows from that monotone convergence: a point of the
closed ball satisfies every bound, while a point outside it violates the bound at
any $n$ with $r+1/n<\norm{c-c_0}_{\mathcal C}$, once $m$ is large. Finally
$B(c_0,r)=\bigcup_{n}\overline{B(c_0,r-1/n)}$ over those $n$ with $r-1/n>0$, so
every open ball, hence every open set, lies in $\Sigma_{\mathrm{cyl}}$.
\end{proof}

The cylinders themselves form a $\pi$-system, the intersection of two of them
being the cylinder over the combined tuple, and because $\Z^2$ is countable that
family generates $\Sigma_{\mathrm{cyl}}$; \cref{thm:cylinder-borel} therefore
exhibits a $\pi$-system generating the Borel sigma algebra, which is what the
determination theorem of \cref{sec:identification} consumes through the extension
argument of Dynkin.

\subsection{The lift, and the topology of convergence}

\Cref{sec:lattice} already exhibits a total carrier-valued representative of the
physical lattice law; this section states it once as a lemma, so that no later
section has to hedge about it.

\begin{lemma}[The finite-source lift is a carrier-valued representative]
\label{lem:lift-representative}
Let $N\ge1$ and let $\mu^{\mathrm{GJ}}_N$ be the physical lattice Gaussian law of
\cref{def:gj-measure}. Then the finite-source Fourier lift
$\mathcal L_N:\R^{I_N}\to\C^{\Z^2}$ of \cref{def:finite-source-lift} takes
values in $\mathcal C$ for \emph{every} configuration, it is Borel as a map into
$\mathcal C$, and its entries are
\begin{equation}\label{eq:lift-coordinate-identity}
  \bigl(\mathcal L_N(\omega)\bigr)_k
  =\sqrt{w_3(k)}\;a^{2}\sum_{x\in I_N}\omega(x)\,e^{-2\pi i\,k\cdot(ax)/L}
  \qquad(\omega\in\R^{I_N},\ k\in\Z^2).
\end{equation}
Consequently $\mu_N=(\mathcal L_N)_*\mu^{\mathrm{GJ}}_N$ is a Borel probability
measure on $\mathcal C$ that carries the law of the lattice field, and no
exceptional set is involved.
\end{lemma}

\begin{proof}
Three facts, each established before this point. \emph{Range and finiteness:} by
\eqref{eq:lift}, $\mathcal L_N\omega$ is a finite real-linear combination of the
coordinate evaluations $\omega\mapsto\omega(x)$, and
$\norm{\mathcal L_N\omega}_{\mathcal C}^{2}$ is finite for every $\omega$ by the
weight sum computed in \cref{sec:lattice}, where it is bounded by
$a^4\norm{\omega}_1^2\bigl(1+\sum_{p\ne0}\norm{p}_\infty^{-6}\bigr)$.
\emph{Measurability:} each coordinate $c\mapsto c_k$ is continuous
(\cref{lem:coordinate-embedding}), hence Borel, and those coordinates generate
$\mathcal B(\mathcal C)$ (\cref{thm:cylinder-borel}); a map into a separable
metric space whose coordinates are Borel is Borel. \emph{The entry formula} is
\eqref{eq:lift}, which carries the factor $\sqrt{w_3}$ of the stored
representation; the entries are therefore the stored ones, not the raw
coefficients, and \cref{rem:two-families} is the dictionary between the two. The
last sentence is the definition of $\mu_N$ in \cref{eq:finite-source-law}
together with those three facts.
\end{proof}

\begin{remark}[What the formalisation does with the lift]\label{rem:lift-development}
The development of \cref{sec:lean-overview} does not use $\mathcal L_N$ for this
purpose. Its carrier-valued lift is a \emph{total} function, defined
coordinatewise (the series appears only in the $\ell^2$ side condition that the
completion imposes, not in the definition), together with the convention that the
value is $0$ off the convergence set. Its measurability is carried as one
conjunct of an input of the formal development rather than derived there; the
other conjunct asserts that this chosen function agrees with the raw weighted
lift almost everywhere, so that input does speak to the size of the exceptional
set. Three things follow, and all are about the formalisation.
First, the input is a statement about one chosen total function, not about the
existence of a carrier-valued representative, which
\cref{lem:lift-representative} settles outright.
Second, what the input supplies is a.e.-measurability and a.e.-agreement,
whereas the lemma supplies a range statement for every $\omega$ and genuine
Borel measurability; a reproducer reading only the input should not read it as
the lemma in other notation.
Third, nothing here turns on that difference on the finite-source route: the
development itself proves, with no input, that the composition of the total lift
with the lattice embedding is measurable, and that its pushforward coincides
with the finite-source pushforward. A reader who compares the two should read
\cref{thm:main} as a statement about $\mu_N$, which is what it is.
\end{remark}


\begin{remark}[The topology used for convergence]\label{rem:topology}
Convergence below is weak convergence of Borel probability measures on the
carrier \cite{billingsley1999}: $\mu_n\Rightarrow\mu$ means
$\int\varphi\,\dd\mu_n\to\int\varphi\,\dd\mu$ for every bounded continuous
$\varphi:\mathcal C\to\R$. This is the topology \cref{thm:separable} makes
available. Three things are deliberately avoided: a weak-star topology on a dual,
since the measures are Borel laws on $\mathcal C$ itself and $\mathcal C$ is
never treated as the dual of a test space; a metric on the space of probability
measures, although a Prokhorov metric exists; and a topology on a space of test
functions, since what is paired with the measures is a bounded continuous
function on the carrier, never a distribution. The operational criterion used
later is not that convergence of finitely many coordinate marginals implies weak
convergence, which is false without a tightness hypothesis, but that a weak
limit, whose existence comes from tightness together with \cref{ext:prokhorov},
is determined by its finite coordinate marginals: continuity of the cylinder maps
transfers the marginal limits to any limit point, and \cref{thm:cylinder-borel}
supplies the sigma algebra equality the determination argument needs. Tightness
itself is the content of \cref{sec:tightness}, not of this section.
\end{remark}
